\section{Repozytoria fazy 02 i 03}\label{sec:faza02}

Trzy repozytoria przygotowują komponenty Second Key, które według planów pracy powstaną po
fazie 01: \path{letsgolegacy.secondkey-capture} (C1 od fazy 02),
\path{letsgolegacy.secondkey-sandbox} (C12) i~\path{letsgolegacy.secondkey-portal} (C11).
W~migawkach (commity \repopath{e56068e}, \repopath{75042c6} i~\repopath{e73546f}
z~2026-09-25) każde z~nich zawiera tylko \path{README.md}, \path{docs/WORKPLAN.md},
\path{LICENSE} i~\path{.gitignore}: nie ma kodu, workflowów CI, testów ani ADR, a~każda
pozycja planu ma status „planned” (jedyny wyjątek to decyzja D7 w~sandboxie, „proposed”).
Rozdział opisuje więc cel, zakres, zaplanowane prace i~ich status oraz zależności od fazy 01
-- wyłącznie tak, jak podają je źródła.

\finding{Portal (C11) jest w~tym rozdziale, bo jest trzecim z~repozytoriów planowanych po
fazie 01, ale jego README i~WORKPLAN przypisują całą jego pracę do fazy 03, a~nie 02.}

\zrodla{\path{letsgolegacy.secondkey-capture/README.md},
\path{letsgolegacy.secondkey-sandbox/README.md},
\path{letsgolegacy.secondkey-portal/README.md}}

\subsection{Umiejscowienie i~bramka fazy 01}\label{sec:faza02-bramka}

Tabela „Repositories” w~README rdzenia przypisuje komponenty do repozytoriów:

\begin{tabularx}{\linewidth}{@{}>{\raggedright\arraybackslash}p{6.6cm} >{\raggedright\arraybackslash}X@{}}
\toprule
Repozytorium & Zawiera \\
\midrule
\path{letsgolegacy.secondkey} (rdzeń) & C0 CLI, C1 capture HTTP (faza 01), C2, C3, C6, C7, C8, C9, C10 \\
\path{letsgolegacy.portcullis} & C5 -- bramka \\
\path{letsgolegacy.secondkey-standards} & C4 -- pakiet standardów \\
\path{letsgolegacy.secondkey-capture} & C1 od fazy 02 (SQL, kolejki, PII) \\
\path{letsgolegacy.secondkey-sandbox} & C12 -- pakowanie wewnątrz granicy klienta \\
\path{letsgolegacy.secondkey-portal} & C11 -- portal dowodów \\
\path{letsgolegacy.secondkey-nopcommerce-bench} & Publiczny okaz: nopCommerce 3.x \\
\bottomrule
\end{tabularx}

\textbf{Bramka fazy 01.} Plan pracy rdzenia nazywa fazę 01 „demo” z~bramką: publiczny
evidence pack (pakiet dowodowy) na nopCommerce 3.x; sekcja „Phase 02 and later” zaczyna się od zastrzeżenia,
że faza 02 startuje dopiero po bramce fazy 01. Plan pracy benchu wskazuje swój ticket P9 jako
bramkę fazy: publikację pakietu dowodowego, kontraktu i~liczb (scenariusze, różnice,
regresje, zabite mutanty); w~migawce P9 ma status „planned” (rozdział~\ref{sec:bench}).
Plany trzech repozytoriów fazy 02 same nie powołują się na tę bramkę; wynika ona z~planów
rdzenia i~benchu.

Z~planu rdzenia wynikają też zależności tych repozytoriów od fazy 01 i~od siebie
nawzajem:

\begin{itemize}
  \item \textbf{Artefakty.} \repopath{*.skcap} (zdarzenia JSONL korelowane identyfikatorem
        żądania) produkuje C1 capture, a~czytają C2 miner i~C6 replay; \repopath{verdict.json}
        czytają C10 i~C11 portal; \repopath{evidence.dsse} (in-toto/DSSE, podpisywany przez
        C10 od fazy 03) czytają C11 portal i~audytor (rozdział~\ref{sec:formaty}).
  \item \textbf{Przeniesienie C1.} W~tabeli fazy 02 rdzenia pozycja „C1 SQL/queues/PII” ma
        opis: przenosi się do \path{letsgolegacy.secondkey-capture} (strona Windows/IIS,
        własny cykl wydań), faza 02. Capture HTTP fazy 01 (ticket S9 rdzenia: proxy nagrywające
        YARP zapisujące \repopath{*.skcap}, tylko ruch przychodzący) jest zrobiony.
  \item \textbf{Zasady projektowe rdzenia}, które te repozytoria realizują: wszystko działa
        wewnątrz granicy klienta, domyślnie bez ruchu wychodzącego (sandbox); tylko
        technologie, które odbiorca już utrzymuje -- .NET, SQL Server, Entra ID, GitHub albo
        Azure DevOps, Kubernetes (sandbox, portal); system legacy jest nietykalny -- capture nie
        wymaga zmiany kodu ani rekompilacji (capture); brak własnej usługi tożsamości, używany
        jest IdP klienta (portal); każdy krok działa z~CLI w~pipeline, a~portal jest oknem, nie
        interfejsem (portal).
  \item \textbf{Znany trudny problem C1} według rdzenia: przypisanie instrukcji SQL do żądania
        HTTP bez zmian w~kodzie jest heurystyką z~jawnie raportowanym marginesem niepewności.
\end{itemize}

\zrodla{\path{letsgolegacy.secondkey/README.md},
\path{letsgolegacy.secondkey/docs/WORKPLAN.md},
\path{letsgolegacy.secondkey-nopcommerce-bench/docs/WORKPLAN.md}}

\subsection{secondkey-capture (C1 od fazy 02)}\label{sec:faza02-capture}

\textbf{Cel.} Nagrywanie tego, jak system legacy naprawdę zachowuje się na swoich granicach --
HTTP, SQL z~parametrami, kolejki -- bez zmiany jego kodu, z~danymi osobowymi (PII)
tokenizowanymi wewnątrz sieci klienta. README wymienia źródła: HTTP, sesje SQL Server
Extended Events i~kolejki.

\textbf{Dlaczego osobne repozytorium.} Capture fazy 01 obsługuje tylko HTTP i~mieszka
w~repozytorium rdzenia. Od fazy 02 capture mieszka tutaj, bo działa po stronie klienta, na
Windows/IIS, i~ma własny cykl wydań.

\textbf{Format wyjścia.} Wynikiem jest format \repopath{*.skcap} zdefiniowany przez schemat
w~repozytorium rdzenia (rozdział~\ref{sec:formaty}); to repozytorium nigdy nie zmienia tego
formatu samodzielnie.

\begin{xltabular}{\linewidth}{@{}>{\raggedright\arraybackslash}p{5.0cm} >{\raggedright\arraybackslash}X >{\raggedright\arraybackslash}p{1.6cm}@{}}
\toprule
Praca (faza 02 -- piloty) & Zrobione, gdy & Status \\
\midrule
\endfirsthead
\toprule
Praca (faza 02 -- piloty) & Zrobione, gdy & Status \\
\midrule
\endhead
Proxy nagrywające YARP, przychodzące i~wychodzące (\repopath{defaultProxy}) & Wymiany (exchange) przychodzące i~wychodzące trafiają do jednego pliku \repopath{.skcap}, skorelowane identyfikatorem żądania & planned \\
Sesja SQL Server Extended Events (\repopath{rpc\_completed}, \path{sql_batch_completed}) & Instrukcje z~parametrami są nagrywane bez dotykania aplikacji & planned \\
Heurystyka korelacji HTTP ↔ SQL (okno czasowe, sesja połączenia) & Każda przypisana instrukcja niesie raportowaną pewność; raportowany jest udział instrukcji nieprzypisanych & planned \\
Kopia kolejek na poziomie brokera & Wiadomości są nagrywane bez zmian w~aplikacji & planned \\
Tokenizacja PII przez sidecar Presidio & Nagrania opuszczające host capture zawierają tokeny, a~nie dane osobowe & planned \\
\bottomrule
\end{xltabular}

Faza 04: usługi Java (modernize-java) -- planned.

\textbf{Odrzucony wariant.} Wczesny plan budowy capture na Keploy porzucono: eBPF działa tylko
na Linuksie, a~legacy .NET Framework działa na Windows/IIS.

\textbf{Pliki w~repozytorium.} \path{.gitignore} (wspólny wzorzec z~sandboxem i~portalem)
ignoruje wyniki builda, pliki IDE, wyniki testów i~testów mutacyjnych, lokalne sekrety
(\path{.env}, \path{secrets.env}; śledzone są tylko pliki przykładowe) oraz wyniki przebiegów
Second Key: nagrania \repopath{*.skcap} i~\repopath{*.skrun} -- z~komentarzem, że nagrania
mogą zawierać dane klienta i~prawdziwych nigdy się nie commituje -- z~wyjątkiem próbek
(\path{samples/}, \path{schemas/samples/}) i~fixture testów, a~także \path{.secondkey/}
i~\path{evidence-out/}. Stan w~README: faza 02; capture HTTP fazy 01 mieszka w~rdzeniu.

\zrodla{\path{letsgolegacy.secondkey-capture/README.md},
\path{letsgolegacy.secondkey-capture/docs/WORKPLAN.md},
\path{letsgolegacy.secondkey-capture/.gitignore},
\path{letsgolegacy.secondkey/docs/WORKPLAN.md}}

\subsection{secondkey-sandbox (C12)}\label{sec:faza02-sandbox}

\textbf{Cel.} „Nothing leaves our network.” Second Key zawsze działa wewnątrz granicy klienta;
to repozytorium pakuje go dla każdego rodzaju takiej granicy: Helm z~pulami węzłów Windows
i~Linux, pojedynczy host Hyper-V, Bicep i~Azure Managed Application dla subskrypcji Azure
klienta oraz paczka dla środowisk odciętych od sieci (air-gapped).

\textbf{Ograniczenie, które kształtuje całość.} Migracje z~.NET Framework wymagają Windows,
więc sandbox niesie oba środowiska: stronę legacy na Windows/IIS, a~zmigrowanego kandydata
i~resztę łańcucha na Linuksie.

\begin{tabularx}{\linewidth}{@{}>{\raggedright\arraybackslash}p{4.4cm} >{\raggedright\arraybackslash}X@{}}
\toprule
Środowisko docelowe & Pakowanie \\
\midrule
Subskrypcja Azure klienta & Azure Managed Application z~Marketplace (zarządzana grupa zasobów, do trybu zablokowanego, w~którym wydawca nie ma dostępu); Bicep \\
Chmura prywatna / on-prem & Chart Helm z~pulami węzłów Windows i~Linux (AKS albo Kubernetes on-prem) \\
Pojedynczy host & Host Windows Server z~Hyper-V: maszyna wirtualna Windows dla strony legacy, maszyna wirtualna Linux dla reszty \\
Air-gapped & Paczka z~obrazami dla własnego rejestru klienta \\
\bottomrule
\end{tabularx}

Według planu tylko demo albo wersja próbna na publicznym kodzie może być hostowana przez
autorów Second Key; każde inne użycie działa w~granicy klienta.

\begin{xltabular}{\linewidth}{@{}>{\raggedright\arraybackslash}p{4.6cm} >{\raggedright\arraybackslash}X >{\raggedright\arraybackslash}p{0.9cm} >{\raggedright\arraybackslash}p{1.6cm}@{}}
\toprule
Praca & Zrobione, gdy & Faza & Status \\
\midrule
\endfirsthead
\toprule
Praca & Zrobione, gdy & Faza & Status \\
\midrule
\endhead
Decyzja o~topologii demo (D7): host Windows z~Hyper-V albo mały klaster z~obiema pulami & ADR zostaje przyjęty & 01 & proposed \\
Chart Helm: pule węzłów Windows i~Linux, obciążenia legacy i~kandydata, worker replay, domyślnie bez ruchu wychodzącego (NetworkPolicy) & \texttt{helm lint} i~\texttt{helm template} przechodzą w~CI; test dymny oparty na kind uruchamia stronę Linux & 02 & planned \\
Wariant jednego hosta (Hyper-V) & Skryptowy build obu maszyn wirtualnych jest udokumentowany i~zlintowany & 02 & planned \\
Bicep dla Azure klienta & \texttt{az bicep build} przechodzi w~CI & 02 & planned \\
Paczka air-gapped & Skrypt tworzy paczkę obrazów i~manifest skrótów (digestów) & 02 & planned \\
Pakiet Azure Managed Application & \path{createUiDefinition} i~\repopath{mainTemplate} przechodzą walidację & 03 & planned \\
\bottomrule
\end{xltabular}

Plan ma jeszcze jedną pozycję fazy 02 (planned), nietechniczną, dotyczącą licencji Windows
Server i~SQL Server w~sandboxie; ten dokument jej nie opisuje.

\textbf{Zależność od fazy 01: decyzja D7.} Jedyna pozycja sandboxu przypisana do fazy 01 to
decyzja o~topologii demo, ze statusem „proposed” i~warunkiem „ADR zostaje przyjęty”. Bench
zapisuje proponowaną decyzję D7 w~\path{docs/TOPOLOGY.md} (status: proposed): jeden host
Windows uruchamia stronę legacy (nopCommerce 3.90 na IIS z~SQL Server), a~Linux -- resztę
łańcucha; tylko capture i~replay muszą sięgać do działającego sklepu przez HTTP, wszystko
dalej to „plik na wejściu, plik na wyjściu” (rozdział~\ref{sec:bench}).

Stan w~README: faza 02.

\zrodla{\path{letsgolegacy.secondkey-sandbox/README.md},
\path{letsgolegacy.secondkey-sandbox/docs/WORKPLAN.md},
\path{letsgolegacy.secondkey-nopcommerce-bench/docs/TOPOLOGY.md},
\path{letsgolegacy.secondkey/docs/WORKPLAN.md}}

\subsection{secondkey-portal (C11)}\label{sec:faza02-portal}

\textbf{Cel.} Okno, przez które zagląda się głównie do odczytu, na pakiety dowodowe dla osób zatwierdzających zmiany -- ryzyko, zatwierdzanie zmian, audyt wewnętrzny --
logujących się przez własnego dostawcę tożsamości (IdP) klienta. Uzasadnienie z~planu:
osoby od ryzyka i~audytu nie czytają SARIF ani YAML. Portal jest oknem na dowody, a~nie
interfejsem: każdy krok nadal działa z~CLI w~pipeline.

\textbf{Faza.} README podaje status „phase 03”, a~wszystkie pozycje planu mają fazę 03.

\textbf{Reguły z~planu.}
\begin{itemize}
  \item \textbf{Brak własnej usługi tożsamości.} Produkcja loguje przez IdP klienta po OIDC
        (Entra ID, ADFS). \repopath{authservice} jest używany tylko dla wersji próbnej
        hostowanej samodzielnie.
  \item Role: \emph{engineer}, \emph{reviewer}, \emph{risk} (tylko odczyt).
  \item Każda decyzja recenzenta trafia do dziennika audytu, do którego można tylko
        dopisywać; te decyzje są też etykietami, które kalibrują triage doradczy C8.
\end{itemize}

\begin{xltabular}{\linewidth}{@{}>{\raggedright\arraybackslash}p{4.8cm} >{\raggedright\arraybackslash}X >{\raggedright\arraybackslash}p{0.9cm} >{\raggedright\arraybackslash}p{1.6cm}@{}}
\toprule
Praca & Zrobione, gdy & Faza & Status \\
\midrule
\endfirsthead
\toprule
Praca & Zrobione, gdy & Faza & Status \\
\midrule
\endhead
Powłoka ASP.NET Core + Blazor z~OIDC do konfigurowalnego IdP & Logowanie działa z~testowym IdP w~CI & 03 & planned \\
Metadane przebiegów w~Postgres (przebiegi, podsumowania werdyktów, skróty pakietów) & Import pakietu z~CLI rdzenia pokazuje przebieg & 03 & planned \\
Widoki dowodów: podsumowanie, per klauzula, różnice per scenariusz & Widoki renderują pakiet-fixture & 03 & planned \\
Decyzje recenzentów z~dziennikiem audytu & Każda decyzja zapisuje kto, kiedy, co i~dlaczego; dziennik jest tylko do dopisywania & 03 & planned \\
Eksport etykiet do kalibracji C8 & Decyzje eksportują się jako oznaczone przykłady & 03 & planned \\
Tenancy wersji próbnej przez \repopath{authservice} & Logowanie w~wersji próbnej działa; konfiguracja produkcyjna od niego nie zależy & 03 & planned \\
\bottomrule
\end{xltabular}

\textbf{Zależności.} Portal czyta artefakty rdzenia: \repopath{verdict.json} i~pakiet
dowodowy z~CLI (w~fazie 01 niepodpisany pakiet C10, ticket S14 rdzenia -- zrobiony), a~od
fazy 03 \repopath{evidence.dsse} (podpisywanie C10 jest w~planie rdzenia pozycją fazy 03).
Etykiety z~decyzji recenzentów zasilają C8 -- triage doradczy, który według planu rdzenia
(fazy 02--03) jest portem kalibracji z~judge-worker (kappa Cohena na etykietach recenzentów,
bez liczby przy małej próbie) i~nigdy nie zmienia werdyktu.

\zrodla{\path{letsgolegacy.secondkey-portal/README.md},
\path{letsgolegacy.secondkey-portal/docs/WORKPLAN.md},
\path{letsgolegacy.secondkey-portal/.gitignore},
\path{letsgolegacy.secondkey/docs/WORKPLAN.md}}

\subsection{Czego źródła nie określają}\label{sec:faza02-luki}

\begin{itemize}
  \item Pozycje planów capture, sandboxu i~portalu nie mają identyfikatorów ticketów (poza D7)
        ani dat i~nie określają kolejności realizacji.
  \item Repozytoria nie mają kodu, CI, testów ani ADR, więc nie ma interfejsów (poleceń, opcji,
        plików konfiguracyjnych) do opisania.
  \item Plan capture nie wskazuje brokera kolejek, formatu tokenów PII ani sposobu wdrożenia
        sidecara Presidio; plan sandboxu -- obrazów, rejestru ani wersji Kubernetes; plan
        portalu -- schematu bazy Postgres ani formatu eksportu etykiet.
\end{itemize}

\zrodla{\path{letsgolegacy.secondkey-capture/docs/WORKPLAN.md},
\path{letsgolegacy.secondkey-sandbox/docs/WORKPLAN.md},
\path{letsgolegacy.secondkey-portal/docs/WORKPLAN.md}}
